% A04. General SAP-free original orbital parity, full-lattice Hasse covers,
% and relative-K exact closure separation. Moved here from later chapters;
% theorem labels retained for downstream references. 20 October 2026.
\section{Finite relative indices and separation of exact equivalences}
\label{sec:foundations-relative-parity}

We use the original, equality-inclusive represented exact algebra of
\cref{lem:exact-product}. Throughout this section $M$ is one actual
transitive homogeneous complete finite symmetric binary source with
semantic homomorphic CCO; \emph{SAP is not assumed}. All invariant binary
relations commute under represented composition: their products are
invariant and symmetric, and taking converses reverses the factors.
Write $\operatorname{ExEq}(M)$ for the full exact equivalences; ordinary
invariant equivalences need not be exact. Every relation called $K$ below
is an actual equivalence on the original real sort, not a formal identity
of a quotient palette.

\begin{lemma}[Finite-relative orbital parity without SAP]
\label{lem:ep-ap-orbital-parity}
Let $E<F$ be arbitrary invariant equivalences on $M$ with finite
relative index $n=[F:E]$. For every old diversity atom
$c\subseteq F\setminus E$, its graph is bipartite. Its even- and
odd-walk relations $Y_c,O_c$ are represented by stabilized positive
atomic-word profiles; $Y_c$ is a full exact equivalence, and
\[
 Y_c<X_c=Y_c\cup O_c\le F,
 \qquad [X_c:Y_c]=2.
\]
\end{lemma}
\begin{proof}
Suppose the $c$-graph contains an odd cycle of length $\ell$.
Choose a connected finite graph $G$ of girth greater than $\ell$
and chromatic number greater than $n$, by the Erd\H{o}s
high-girth/high-chromatic theorem. Label all its edges by $c$.
Every specified simple cycle of $G$ has a colour-preserving
homomorphism into the old $c$-graph: an even cycle folds onto one
$c$-edge, while any odd cycle (necessarily of length greater than
$\ell$) traverses the old odd $\ell$-cycle with two-step
backtracks. Semantic CCO therefore yields a homomorphism $G\to M$.
Since $G$ is connected and $c\subseteq F$, its image lies in a
single $F$-class. Since $c\cap E=\varnothing$, adjacent vertices
map to different $E$-classes. Their $n$ classes give a proper
$n$-colouring of $G$, a contradiction.

By transitivity every real vertex has a $c$-neighbour. Inserting
one two-edge backtrack in a represented $c$-walk shows
$c^j\subseteq c^{j+2}$. Finiteness of the old orbital palette
makes the even and odd power chains stabilize. The stable even
word profile is exactly the even-walk relation $Y_c$, and the
stable odd profile is $O_c$; repeated intermediate real vertices
are allowed by the represented multiplication convention.
Every nonempty connected component of a bipartite $c$-graph
has exactly two parity classes. Hence $Y_c$ is an equivalence,
$X_c=Y_c\cup O_c$ is its connectivity equivalence, and their
relative index is two. The stabilized even positive word is
full exact, without an injectivity or SAP assumption.
\end{proof}

\begin{theorem}[Finite Hasse indices in the full invariant lattice]
\label{thm:full-relative-Hasse}
Every finite-relative-index cover $E\prec F$ in the \emph{full}
lattice of invariant equivalences of $M$ has $[F:E]=2$.
Consequently every finite relative index $[F:E]$ (not necessarily
covering) is a power of two.
\end{theorem}
\begin{proof}
There are finitely many invariant equivalences, since they are
unions of finitely many orbitals. They permute under represented
composition, so the lattice is modular. Induct on the height of
the upper endpoint $F$. Choose an old atom
$c\subseteq F\setminus E$ and apply
\cref{lem:ep-ap-orbital-parity} to obtain
$Y_c\prec X_c\le F$ with $[X_c:Y_c]=2$ and $X_c\not\le E$.
Since $E\prec F$, we have $E\vee X_c=F$. Put $K=E\wedge X_c$.
Dedekind interval transposition for permuting equivalences gives
\[
 K\prec X_c,\qquad [X_c:K]=[F:E].
\]
If $K=Y_c$ the index is two. Otherwise $K,Y_c$ are distinct
coatoms of $X_c$, so $K\vee Y_c=X_c$. Put $L=K\wedge Y_c$.
A second transposition gives
\[
 L\prec Y_c,\qquad [Y_c:L]=[X_c:K]=[F:E].
\]
Since $Y_c<X_c\le F$, the upper endpoint $Y_c$ has smaller
height than $F$, and the induction hypothesis makes this index
two. For any finite $[F:E]$, multiply the indices along a
saturated chain between $E$ and $F$.
\end{proof}

\begin{lemma}[Relative exact closure separation]
\label{lem:ep-relative-K-closure}
Let $K$ be a full exact invariant equivalence on $M$. Let
$A,B$ be nonempty symmetric full exact $K$-bisaturated relations
containing $K$, and suppose $A\cap B=K$. Their stabilized
represented equivalence closures satisfy
\[
 \boxed{\langle A\rangle\cap\langle B\rangle=K.}
 \tag{S2}\label{eq:ep-relative-K-closure}
\]
Here $\langle A\rangle$ denotes the stable power of $A$;
this convention includes $K$ as the zeroth power.
\end{lemma}
\begin{proof}
Set $A^{[0]}=B^{[0]}=K$ and use ordinary represented powers
for positive exponents. Since $K\subseteq A,B$ and both relations
are $K$-bisaturated, the powers form increasing chains of
symmetric full exact $K$-bisaturated relations. They stabilize
to full exact equivalences. Suppose an old atom outside $K$
belongs to both stable closures. Choose the least $n\ge1$ for
which $(B^{[n]}\cap\langle A\rangle)\setminus K$ is nonempty,
choose an old atom $c$ in that difference, and the least
$m\ge1$ such that $c\in A^{[m]}$.
We claim
\[
 c\notin A^{[m-1]};B^{[n-1]}.
 \tag{S3}\label{eq:ep-minimal-c}
\]
Otherwise write $c\in a;b$ for atoms
$a\in A^{[m-1]}$, $b\in B^{[n-1]}$. Peircean transfer gives
$b\in a^\smile;c\subseteq\langle A\rangle$. If $b\notin K$,
this contradicts minimality of $n$; if $b\in K$, then
$K$-saturation puts $c\in A^{[m-1]}$, contradicting
minimality of $m$.

Set $T=A^{[m-1]};\{c\};B^{[n-1]}$. Full exact product
closure makes $T$ nonempty exact and $K$-bisaturated. By
Peircean rotation of $c\in A^{[m-1]};A$ and
$c\in B;B^{[n-1]}$, the relation $T$ meets both $A$ and $B$.
If $T\cap K$ were nonempty, invariance and $K$-bisaturation
would give $\Id\subseteq T$, which by symmetry and
Peircean rotation would yield
$c\in A^{[m-1]};B^{[n-1]}$, contrary to the claim.
Thus $A,B,T$ are nonempty exact domains meeting pairwise,
but $A\cap B\cap T=\varnothing$, contrary to full exact
$2$-Helly, \cref{lem:exact-product}.
\end{proof}

\begin{corollary}[Separation at the old logical identity]
\label{cor:foundations-identity-separation}
For full exact symmetric $A,B$ containing $\Id$ and with
$A\cap B=\Id$, their exact equivalence closures intersect in
$\Id$. This is the case $K=\Id$ of
\cref{lem:ep-relative-K-closure}.
\end{corollary}

\begin{theorem}[Maximal exact parity gate and its source involution]
\label{thm:foundational-maximal-parity-gate}
Let $Y<X$ be actual invariant same-sort equivalences in the
original homogeneous symmetric AP+homomorphic-CCO source,
with $[X:Y]=2$. Assume $Y$ and the \emph{odd relation}
$O=X\setminus Y$ are full exact. The equivalence $X$
itself need \emph{not} be full exact. Choose an
inclusion-maximal full exact equivalence $K$ satisfying
\[
             Y\le K,\qquad K\cap X=Y.
\]
Then, writing $F=K\vee X$ and $t=K;O;K$,
\[
             [F:K]=2,\qquad
             F=K\mathbin{\dot\cup}t,\qquad
             t;t=K.
\]
The relation $t$ is nonempty, full exact and $K$-bisaturated;
its image $\bar t=t/K$ is the graph of an equivariant
fixed-point-free involution of $M/K$.
More strongly, \emph{every original full exact
equivalence} $Z>K$ contains the entire
old relation $t$. In particular every full
exact quotient equivalence
$\bar Z>\Id_{M/K}$ contains $\bar t$.
The first assertion does not require any
converse exact-descent theorem.
Thus the original source object with kernel $K$ has the
distinguished nonidentity exact isotropy arrow required by
\cref{thm:foundational-rigid-character}.
No SAP is required.
\end{theorem}
\begin{proof}
Existence of maximal $K$ follows from finiteness
of the invariant-equivalence lattice; the initial
choice $K=Y$ is admissible and full exact.
All invariant equivalences permute under represented
composition because same-sort orbital relations
are symmetric. Dedekind transposition gives
$K\prec K\vee X$ and $[F:K]=[X:Y]=2$.
Since $Y\le K$ and $X=Y\cup O$,
\[
 F=KX=K\cup KO=K\cup KOK=K\cup t.
\]
The unions are disjoint: an edge in $t\cap K$
would force an $O$-edge inside one $K$-class,
contradicting $O\cap K=\varnothing$.
Full exactness of $K,O$ and exact convolution
make $t$ full exact. Because every $F$-class
contains exactly two $K$-classes, $t$ is
the unique matching between those classes,
so it is a fixed-point-free involutive graph
and $t^2=K$ as represented relations.

First take \emph{any original} full exact
invariant equivalence $Z>K$. If
$Z\cap X$ contains an old atom of $O$,
then $Z\cap X$ is an equivalence strictly
above $Y$, and since $[X:Y]=2$ it equals
$X$. Thus $O\subseteq Z$ and
$t=K;O;K\subseteq Z$. Otherwise
$Z\cap X=Y$, which contradicts
maximality of $K$. Hence the entire
relation $t$ belongs to every original
full exact equivalence strictly above $K$.

Now take any nontrivial full exact quotient
equivalence $\bar Z$ on $M/K$. Its old
inverse image $Z$ is original full exact
by the typed exact quotient word-pullback
\cref{lem:typed-simultaneous-blocks}, and
$Z>K$. The preceding \emph{old-level}
claim gives $t\subseteq Z$, so
$\bar t\subseteq\bar Z$. Importantly,
no step requires descent of an arbitrarily
transported old exact relation to the
quotient. The relation $t$ is a bridge by
\cref{def:foundational-exact-bridge} with
source and target kernel $K$.
\end{proof}
